\begin{proposition}[6.13.2]
\label{IV.2.s6.13.2-fr}
Sean $X$ un preesquema localmente noetheriano reducido y $X'$ su
normalización \textup{(\textbf{II}, 6.3.8)}. Si el morfismo canónico $f:X'\to X$ es finito,
$\operatorname{Nor}(X)$ es abierto en $X$.
\end{proposition}

\begin{proof}
En efecto, se tiene $X'=\operatorname{Spec}(f_*(\mathcal O_{X'}))$ y
$f_*(\mathcal O_{X'})$ es una $\mathcal O_X$-álgebra \emph{coherente} por
hipótesis (\textbf{II}, 6.1.3). Decir que $X$ es normal en un punto $x$ significa
que el homomorfismo canónico
$\mathcal O_x\to(f_*(\mathcal O_{X'}))_x$ es biyectivo (\textbf{II}, 6.3.4);
pero el conjunto de estos puntos es \emph{abierto}, puesto que $\mathcal O_X$
y $f_*(\mathcal O_{X'})$ son coherentes (\textbf{0}$_{\rm I}$, 5.2.7).
\end{proof}
